\clearpage
\VerdierPhysicalPageAnchor{0235}%
\begingroup\small\setlength{\parskip}{0pt}\setlength{\emergencystretch}{1.5em}
\setlength{\abovedisplayskip}{4pt plus 2pt minus 2pt}%
\setlength{\belowdisplayskip}{4pt plus 2pt minus 2pt}%

\VerdierSemanticAnchor{ch4d-4-2-2-cont}
\noindent where $(p,q)$ and $(p',q')$ are elements of
$\mathcal{F}\ell_1(\widetilde{\mathbb Z})$ such that $p\leq p'$ and
$q\leq q'$, the image in $D(\mathcal A)$ of the canonical morphism in
$\operatorname{comp}(\mathcal A)$:

\VerdierSemanticAnchor{ch4d-4-2-2-3}
\begin{equation}
\psi\bigl((p,q),(p',q')\bigr)\colon
X(p,q)\longrightarrow X(p',q').
\tag{4.2.2.3}
\end{equation}
We have thus defined a functor from
$\mathcal{F}\ell_1(\widetilde{\mathbb Z})$ to $D(\mathcal A)$. For every
element $(p,q,r)$ of
$\mathcal{F}\ell_2(\widetilde{\mathbb Z})$, let

\VerdierSemanticAnchor{ch4d-4-2-2-4}
\begin{equation}
S(p,q,r)\colon
0\longrightarrow X(p,q)\longrightarrow X(p,r)
\longrightarrow X(q,r)\longrightarrow0
\tag{4.2.2.4}
\end{equation}
be the exact sequence in $\operatorname{comp}(\mathcal A)$ whose
morphisms are the canonical morphisms induced by the filtration of $X$.
Finally, for every element $(p,q)$ of
$\mathcal{F}\ell_1(\widetilde{\mathbb Z})$, let

\VerdierSemanticAnchor{ch4d-4-2-2-5}
\begin{equation}
J(p,q)\colon\mathcal X(p,q)\longrightarrow X(p,q)
\tag{4.2.2.5}
\end{equation}
be the image in $D(\mathcal A)$ of the morphism in
$\operatorname{comp}(\mathcal A)$

\VerdierSemanticAnchor{ch4d-4-2-2-6}
\begin{equation}
j(p,q)\colon\mathcal X(p,q)\longrightarrow X(p,q),
\tag{4.2.2.6}
\end{equation}
whose components are

\VerdierSemanticAnchor{ch4d-4-2-2-7}
\begin{equation}
\left\{
\begin{aligned}
j(p,q)^i\colon\mathcal X(p,q)^i
 &=X(q)^i\oplus X(p)^{i+1}
   \longrightarrow X(q)^i/X(p)^i=X(p,q)^i,\\
j(p,q)^i&=(\pi^i,0),\\
&\text{where $\pi^i$ is the canonical epimorphism:}\\
\pi^i\colon X(q)^i&\longrightarrow X(q)^i/X(p)^i.
\end{aligned}
\right.
\tag{4.2.2.7}
\end{equation}

\VerdierSemanticAnchor{ch4d-4-2-3}
\medskip\noindent\textbf{Proposition 4.2.3.}\quad
\textit{For every element $(p,q,r)$ of
$\mathcal{F}\ell_2(\widetilde{\mathbb Z})$, denote by}
\[
\delta(p,q,r)\colon X(q,r)\longrightarrow X(p,q)[1]
\]
\noindent\textit{the morphism in $D(\mathcal A)$}

\VerdierSemanticAnchor{ch4d-4-2-3-1}
\begin{equation}
\delta(p,q,r)=\delta\bigl(S(p,q,r)\bigr),
\tag{4.2.3.1}
\end{equation}
\noindent\textit{where $\delta\bigl(S(p,q,r)\bigr)$ is the morphism in
$D(\mathcal A)$ associated with the exact sequence $S(p,q,r)$ (4.2.2.4)
by the canonical $\delta$-functor
$\operatorname{comp}(\mathcal A)\longrightarrow D(\mathcal A)$ (1.3.2).}

\VerdierSemanticAnchor{ch4d-4-2-3-ab}
\noindent\textit{a) The objects $X(p,q)$ (4.2.2.1), the morphisms
$\varphi\bigl((p,q),(p',q')\bigr)$ (4.2.2.2), and the morphisms
$\delta(p,q,r)$ (4.2.3.1) define a spectral object with values in
$D(\mathcal A)$.}

\noindent\textit{b) The morphisms $J(p,q)$ (4.2.2.5) define an
isomorphism from the spectral object
$\bigl(\mathcal X(p,q),\delta(p,q,r)\bigr)$ onto the spectral object
$\bigl(X(p,q),\delta(p,q,r)\bigr)$.}

\VerdierSemanticAnchor{ch4d-4-2-3-proof}
Assertion a) follows immediately from the properties of the canonical
$\delta$-functor
$\operatorname{comp}(\mathcal A)\longrightarrow D(\mathcal A)$ (1.3.1).
For the proof of assertion b), the reader is referred to (1.3.2).
\endgroup
